% Comparison: bit-serial schoolbook vs word-level clmul leaf
% Used in §6.2.1 (ssec:bit-to-word)

\tikzsetnextfilename{ch6-polymul-schoolbook-vs-clmul}
\begin{tikzpicture}[
  box/.style={draw, rounded corners=2pt, minimum width=4.8cm,
    minimum height=2.2cm, align=left, font=\ttfamily\small},
  arrow/.style={->, thick, >=stealth},
  label/.style={font=\sffamily\small\bfseries, above}
]

% Left box: bit-serial schoolbook
\node[box, fill=red!5] (school) at (0,0) {%
  \begin{tabular}{@{}l@{}}
  for i in 0..n-1:\\
  \quad for bit in 0..63:\\
  \quad\quad mask = -(a[i]>>bit \& 1)\\
  \quad\quad for j in 0..n-1:\\
  \quad\quad\quad r[i+j] \^{}= (b[j]<<bit) \& mask\\
  \quad\quad\quad r[i+j+1] \^{}= (b[j]>>inv) \& mask
  \end{tabular}
};
\node[label] at (school.north) {Bit-serial schoolbook};
\node[font=\sffamily\footnotesize, below=2pt] at (school.south)
  {$64 \times n^2$ conditional shift-XOR};

% Right box: clmul leaf
\node[box, fill=green!5] (clmul) at (7.5,0) {%
  \begin{tabular}{@{}l@{}}
  for i in 0..n-1:\\
  \quad for j in 0..n-1 (pairs):\\
  \quad\quad lo = clmul(a[i], b[j])\\
  \quad\quad hi = clmulh(a[i], b[j])\\
  \quad\quad r[i+j] \^{}= lo\\
  \quad\quad r[i+j+1] \^{}= hi
  \end{tabular}
};
\node[label] at (clmul.north) {Word-level \textsf{Zbc} leaf};
\node[font=\sffamily\footnotesize, below=2pt] at (clmul.south)
  {$n^2$ ζεύγη clmul/clmulh};

% Arrow between
\draw[arrow, very thick, color=blue!60!black]
  (school.east) -- node[above, font=\sffamily\small] {$\times$64}
  (clmul.west);

\end{tikzpicture}
